\documentclass[11pt,a4paper]{article}
\usepackage[T1]{fontenc}
\usepackage[utf8]{inputenc}
\usepackage{lmodern}
\usepackage[margin=25mm,headheight=14pt]{geometry}
\usepackage{amsmath,amssymb,amsthm,mathtools,mathrsfs}
\usepackage{microtype,booktabs,array,longtable}
\usepackage{enumitem,fancyhdr,xurl}
\usepackage[hidelinks]{hyperref}
\usepackage[nameinlink,noabbrev]{cleveref}
\setlength{\parskip}{3pt}
\setlength{\parindent}{1.25em}
\setlength{\emergencystretch}{2em}
\allowdisplaybreaks[2]
\pagestyle{fancy}
\fancyhf{}
\fancyhead[L]{\small Uniform Gaussian-binomial crossover}
\fancyhead[R]{\small Research manuscript}
\fancyfoot[C]{\thepage}
\newtheorem{theorem}{Theorem}[section]
\newtheorem{proposition}[theorem]{Proposition}
\newtheorem{lemma}[theorem]{Lemma}
\newtheorem{corollary}[theorem]{Corollary}
\theoremstyle{definition}
\newtheorem{definition}[theorem]{Definition}
\theoremstyle{remark}
\newtheorem{remark}[theorem]{Remark}
\newcommand{\e}{\mathrm e}
\newcommand{\ii}{\mathrm i}
\newcommand{\dd}{\,\mathrm d}
\newcommand{\Li}{\operatorname{Li}}
\newcommand{\Res}{\operatorname{Res}}
\newcommand{\N}{\mathbb N}
\newcommand{\R}{\mathbb R}
\newcommand{\C}{\mathbb C}
\newcommand{\Borel}{\mathcal B}
\newcommand{\Ecal}{\mathcal E}
\newcommand{\Acal}{\mathcal A}
\newcommand{\Jcal}{\mathcal J}
\newcommand{\Dcal}{\mathcal D}
\newcommand{\Pcal}{\mathcal P}
\DeclareMathOperator{\csch}{csch}
\title{\textbf{Uniform Resurgent Crossover for\\Gaussian Binomial Coefficients}\\[5pt]
\large Exact remainders, a sharp modular-visibility threshold,\\and certified inversion}
\author{Research manuscript prepared for Vladimir Reshetnikov}
\date{29 September 2026}
\begin{document}
\maketitle
\begin{abstract}
The transseries companion in the ProveIt repository develops the singular
$q\to1$ crossover of central Gaussian coefficients, but explicitly leaves
uniform matching at the classical endpoint and the finite-size tail's
beyond-all-orders accuracy untreated. We address these questions for the
specified analytic interpolation, rather than for arbitrary interpolations
of the same sequence. A positive double-kernel representation supplies an
exact signed remainder after every Bernoulli truncation. It yields an
all-order expansion uniform for $x\to\infty$ and $0\le hx\le T$, including the
ordinary central-binomial endpoint. We compute the meromorphic Borel
transform, its parity cancellations, and its imaginary-axis residue jump.
The principal quantitative result is the uniform sharp estimate
\[
 \min_{K\ge0}r_K(\tau/x,\tau)
 =\frac{\e^{-2\pi x}}{\pi\sqrt{x}}\bigl(1+O_T(x^{-1})\bigr),
 \qquad 0\le\tau\le T,
\]
with every minimizing index $K=\pi x+O_T(1)$. Comparison with the exact
modular correction $\Ecal(h)\sim\e^{-4\pi^2/h}$ gives a transition at
$\tau=hx=2\pi$, and a resolved boundary layer
$\tau=2\pi-(\log x)/(2x)+s/x$ in which the error-to-signal ratio tends to
$\e^{-s}/\pi$. This is a truncation-accuracy transition, not a Stokes
transition. Finally, a uniform slope bound transfers residual certificates
to inverse errors, and inversion about the exact Borel-summed carrier gives
a convergent modular series with finite divisor-sum coefficients. Full
proofs, reproducible numerical checks, a source audit, and further research
questions are included. General $q$-Pochhammer resurgence is established
prior work; the contribution claimed here is the explicit uniform,
sharp-remainder, and inversion package, not priority for that general theory.
\end{abstract}
\vspace{3pt}
\noindent\textbf{Status.} Conventional mathematical proofs and exact symbolic
checks; no Lean verification or peer review is claimed. Novelty relative to
the inspected repository chapter is identified below. Literature-wide
priority has not been established.
\clearpage
{\small\setlength{\parskip}{0pt}\tableofcontents}
\clearpage

\section{Problem, source boundary, and principal contributions}
\label{sec:context}

\subsection{The precise repository question}
The repository snapshot used here is commit
\texttt{04e06e032dff1966513bfba316966d52db3646b4} of ProveIt. The directly
relevant source is the chapter \emph{The singular transition $q\to1$} in
\emph{Combinatorial Transseries and Their Inverses}; see
\cite{repo,reporeadme}. Its finite-order tail proposition
\texttt{t3:prop:double-A} is uniform only for $\tau\ge\tau_0>0$.
The discussion following \texttt{t3:eq:Ssmall} explicitly says that uniform
matching at $\tau=0$ requires another error analysis. The discussion of
\texttt{t3:eq:modularsectors} and
\texttt{t3:eq:central-double-asympt} separates an exact modular product from
an algebraic remainder and does not claim to have resolved the finite-size
tail's entire beyond-all-orders contribution. Those are the two scoped
questions developed in this paper.

We inspected the companion's relevant source ranges, approximately
lines 4090--4250, and the two group READMEs. The much larger canonical
\emph{Transseries and Inversion} source could not be retrieved through the
available file endpoint. Its README and the group concordance were read,
but no exhaustive nonduplication audit of that entire volume is claimed.
The repository itself carefully distinguishes human mathematical statements
from its narrower Lean crosswalk; this manuscript preserves that distinction.

The crossover already has the correct dilogarithmic leading phase, a
positive phase derivative, formal higher inverse coefficients, and an exact
eta-modular factor in the companion. These are starting points, not results
claimed as discoveries here. The questions are quantitative: can one cross
the singular endpoint uniformly, what is the \emph{smallest actual}
truncation error, and when can the much smaller modular contribution be
resolved by that representation?

\subsection{What is proved, and what is not}
Our main results are as follows; all refer to the explicitly defined
interpolation in \cref{sec:normalization}.
\begin{enumerate}[label=(\roman*),leftmargin=2em]
\item \Cref{thm:remainder,thm:uniform} give an exact positive remainder and
an all-order expansion uniform on the closed crossover range $0\le hx\le T$.
\item \Cref{thm:borel,prop:scaledborel} give the explicit Borel pole divisor,
its central parity cancellation, and its classical endpoint limit.
\item \Cref{thm:optimal,thm:visibility} determine the optimal error including
its leading constant and the logarithmically shifted modular-visibility
boundary layer. The lower bound concerns all ordinary truncations, not only
a suggested least-term index.
\item \Cref{thm:slope,thm:inverse} provide real residual certificates and a
convergent modular inverse series about the exact, nonmodular carrier.
\end{enumerate}
The proofs do not establish resurgence of every inverse formal series,
a general Hahn-field construction, or a solution of the wider median-sum
conjectures in the literature. Nor does a lower bound for ordinary
truncations prohibit a more accurate integral, hyperasymptotic, or exact
product algorithm. An exact positive remainder is itself a way around the
truncation limitation.

\subsection{Relation to established literature}
Costin and Garoufalidis developed resurgence for Euler--Maclaurin
summation \cite{costingaroufalidis}. Garoufalidis and Kashaev give an explicit
meromorphic Borel transform, poles, residues, and Stokes behavior for
Faddeev's quantum dilogarithm \cite{garoufalidiskashaev}. Fantini and Rella
study $q$-Pochhammer summability, quantum modularity, and character-weighted
modular resurgent series; the version reviewed here is dated 1 April 2026
\cite{fantinirella}. Thus neither the use of Borel transforms for
$q$-products nor the existence of their pole lattices is new in itself.

The elementary derivation below is included to fix all signs and
normalizations and to make the central quotient's cancellations available
for the uniform remainder theorem. The narrow contribution advanced for
further expert checking is the assembled theorem with a uniform endpoint,
the sharp $1/(\pi\sqrt{x})$ error constant, its precision boundary layer,
and the explicitly certified inverse construction. The source comparison
identifies a genuine extension of the inspected chapter; it is not a claim
that no equivalent theorem exists elsewhere.

\section{Normalization and the two small scales}
\label{sec:normalization}
For $h>0$ put $Q=\e^{-h}$ and define
\begin{equation}
 (a;Q)_\infty=\prod_{j=0}^{\infty}(1-aQ^j),\qquad
 P(Q)=(Q;Q)_\infty.
 \label{eq:poch}
\end{equation}
Our real interpolation, for $x>0$, is
\begin{equation}
 H_h(x)=\e^{hx^2}
 \frac{(Q^{x+1};Q)_\infty^2}
 {P(Q)(Q^{2x+1};Q)_\infty},\qquad F_h(x)=\log H_h(x).
 \label{eq:H}
\end{equation}
For integer $n\ge0$, it is exactly the central Gaussian coefficient
$\bigl[\begin{smallmatrix}2n\\n\end{smallmatrix}\bigr]_{q}$ with $q=\e^h$:
the usual $q\leftrightarrow q^{-1}$ factor is $q^{n^2}$. All logarithms in
real-variable formulas are real logarithms of positive quantities.

The standard $q$-gamma definition gives
\begin{equation}
 H_h(x)=\e^{hx^2}\frac{\Gamma_Q(2x+1)}{\Gamma_Q(x+1)^2},
 \qquad
 H_0(x):=\frac{\Gamma(2x+1)}{\Gamma(x+1)^2}.
 \label{eq:qgamma}
\end{equation}
The fixed-$x$ limit as $h\downarrow0$ is $H_0(x)$, by the $q$-gamma limit
\cite{dlmfqgamma}. The parameter $q$ here is continuous; the limit $q\to1$
is not a limit through orders of finite fields.

Write $\tau=hx$. The tail and modular correction are
\begin{align}
 A(h,\tau)&=\sum_{m\ge1}\frac{\e^{-m\tau}}{m(\e^{hm}-1)}
          =-\log(Q\e^{-\tau};Q)_\infty,\label{eq:A}\\
 z(h)&=\e^{-4\pi^2/h},\qquad
 \Ecal(h)=-\log P(z(h))
   =\sum_{j\ge1}\sigma_{-1}(j)z(h)^j,\label{eq:modular}\\
 \sigma_{-1}(j)&=\sum_{d\mid j}\frac1d.\notag
\end{align}
We set $\Ecal(0)=0$. Absolute convergence proves the divisor sum by
expanding $-\log(1-z^m)$ and grouping the product of the two indices.
The eta transformation \cite{dlmfeta} yields
\begin{equation}
 P(\e^{-h})=\sqrt{\frac{2\pi}{h}}
 \exp\!\left(-\frac{\pi^2}{6h}+\frac h{24}\right)P(z(h)).
 \label{eq:eta}
\end{equation}
Consequently, exactly,
\begin{equation}
 F_h(x)=\frac{\tau^2+\pi^2/6}{h}
 -\frac12\log\frac{2\pi}{h}-\frac h{24}
 +A(h,2\tau)-2A(h,\tau)+\Ecal(h).
 \label{eq:exactforward}
\end{equation}

The repository's phase and amplitude are
\begin{align}
 S(\tau)&=\tau^2+\frac{\pi^2}{6}
          +\Li_2(\e^{-2\tau})-2\Li_2(\e^{-\tau}),\label{eq:S}\\
 J(\tau)&=\frac12\log(1-\e^{-2\tau})-\log(1-\e^{-\tau})
          =\frac12\log\coth(\tau/2).\label{eq:J}
\end{align}
Differentiating the dilogarithms gives
\begin{equation}
 S'(\tau)=2\log(1+\e^\tau),\qquad
 S''(\tau)=\frac{2\e^\tau}{1+\e^\tau},\qquad S(0)=0.
 \label{eq:Sder}
\end{equation}
The phase $S$ is strictly increasing. The two potentially competing
exponentials in the crossover are $\e^{-2\pi x}$ and
$\e^{-4\pi^2/h}$. The former will arise as a \emph{minimal remainder size};
the latter is an actual, separately specified term in an exact identity.
These interpretations must not be interchanged.

\section{An exact positive remainder}
\label{sec:remainder}

\subsection{From a partial fraction identity to a signed error}
The standard partial fractions for $\coth$ \cite{dlmfcoth} give, for $v>0$,
\begin{equation}
 \frac1{\e^v-1}=\frac1v-\frac12
       +2v\sum_{n\ge1}\frac1{(2\pi n)^2+v^2}.
 \label{eq:coth}
\end{equation}
The positive summands permit Tonelli's theorem in \eqref{eq:A}. Thus
\begin{equation}
 A(h,\tau)=\frac{\Li_2(\e^{-\tau})}{h}
 +\frac12\log(1-\e^{-\tau})
 +2h\sum_{m,n\ge1}\frac{\e^{-m\tau}}{(2\pi n)^2+(mh)^2}.
 \label{eq:Akernel}
\end{equation}
Define
\begin{align}
 w_m(\tau)&=2\e^{-m\tau}-\e^{-2m\tau}>0,\label{eq:weights}\\
 R(h,\tau)&=2h\sum_{m,n\ge1}
        \frac{w_m(\tau)}{(2\pi n)^2+(mh)^2},\label{eq:R}\\
 a_k(\tau)&=\frac{|B_{2k}|}{(2k)!}
  \left(2\Li_{2-2k}(\e^{-\tau})-\Li_{2-2k}(\e^{-2\tau})\right)>0.
 \label{eq:ak}
\end{align}
Here $v/(\e^v-1)=\sum_{j\ge0}B_jv^j/j!$ fixes the Bernoulli convention.
The elementary equality
$|B_{2k}|/(2k)!=2\zeta(2k)/(2\pi)^{2k}$ can also be obtained by comparing
Taylor coefficients in \eqref{eq:coth}.

\begin{theorem}[Exact signed remainder and enveloping bounds]
\label{thm:remainder}
For every $h,\tau>0$ and integer $K\ge0$, put
\begin{equation}
 r_K(h,\tau)=2h^{2K+1}\sum_{m,n\ge1}
 \frac{w_m(\tau)m^{2K}}
 {(2\pi n)^{2K}\bigl((2\pi n)^2+(mh)^2\bigr)}.
 \label{eq:rK}
\end{equation}
Then
\begin{align}
 R(h,\tau)&=\sum_{k=1}^K(-1)^{k-1}a_k(\tau)h^{2k-1}
                      +(-1)^K r_K(h,\tau),\label{eq:exactR}\\
 0&<r_K(h,\tau)<a_{K+1}(\tau)h^{2K+1}.\label{eq:nextterm}
\end{align}
In particular, adjacent partial sums bracket $R$. The exact forward
identity is
\begin{align}
 F_h(x)={}&\frac{S(\tau)}h-\frac12\log\frac{2\pi}{h}
 +J(\tau)-\frac h{24}\notag\\
 &+\sum_{k=1}^K(-1)^ka_k(\tau)h^{2k-1}
 +\Ecal(h)+(-1)^{K+1}r_K(h,\tau).
 \label{eq:forwardK}
\end{align}
\end{theorem}
\begin{proof}
For $a,b>0$, finite geometric division gives
\[
 \frac1{a^2+b^2}
 =\sum_{j=0}^{K-1}\frac{(-1)^jb^{2j}}{a^{2j+2}}
  +\frac{(-1)^Kb^{2K}}{a^{2K}(a^2+b^2)}.
\]
Apply this with $a=2\pi n$, $b=mh$, and multiply by $2h w_m$.
Each coefficient sum is absolutely convergent: its $m$-part is an
exponentially weighted polynomial, and its $n$-part is a zeta value of
order at least two. Summing gives \eqref{eq:exactR} and \eqref{eq:rK}.
Replacing the final denominator $a^2+b^2$ by $a^2$ proves the strict
upper bound \eqref{eq:nextterm}. Positivity is immediate. Substituting
\eqref{eq:Akernel} into \eqref{eq:exactforward} proves
\eqref{eq:forwardK} with the displayed sign.
\end{proof}

\subsection{Why this is stronger than a formal Bernoulli expansion}
A finite-order $O(h^{2K+1})$ estimate does not identify the sign, is not
necessarily uniform as $\tau\downarrow0$, and cannot justify choosing
$K$ of order $1/h$. In contrast, \eqref{eq:rK} is an exact identity for
\emph{every} $K$, with no restriction relating $K,h,\tau$. We may therefore
study the simultaneous limit $K\asymp x$ directly rather than substitute a
moving $K$ into a theorem proved only for fixed $K$.

There is an additional structural advantage. As a function of $z=h^2$,
$R(h,\tau)/h$ is a positive Stieltjes-type sum of kernels
$1/(a+bz)$. The resulting remainder sequence is a moment sequence.

\begin{lemma}[Strict log-convexity in the truncation index]
\label[lemma]{lem:logconvex}
For fixed $h,\tau>0$,
\begin{equation}
 r_K^2<r_{K-1}r_{K+1}\quad(K\ge1).
 \label{eq:logconvex}
\end{equation}
The quotients $r_{K+1}/r_K$ are strictly increasing, and $r_K\to\infty$
as $K\to\infty$. In particular, a minimum over $K\ge0$ exists, with at
most two adjacent minimizing indices.
\end{lemma}
\begin{proof}
Let
\[
 \nu_{m,n}=\frac{2h w_m}{(2\pi n)^2+(mh)^2},\qquad
 t_{m,n}=\left(\frac{mh}{2\pi n}\right)^2.
\]
Equation \eqref{eq:rK} says $r_K=\sum\nu_{m,n}t_{m,n}^K$.
Cauchy--Schwarz applied to the two powers $K-1$ and $K+1$ gives
\eqref{eq:logconvex}; equality would force all $t_{m,n}$ with positive
weight to coincide, which they do not. A single summand with $t_{m,n}>1$
proves divergence. Increasing successive quotients prove the final claim.
\end{proof}

\section{Uniform closure of the classical endpoint}
\label{sec:uniform}
Define the normalized coefficient functions
\begin{align}
 \Acal(\tau)&=\frac{S(\tau)}\tau, &\Acal(0)&=2\log2,\label{eq:Ac}\\
 \Jcal(\tau)&=\frac12\log\!\left(\frac\tau2\coth\frac\tau2\right),
 &\Jcal(0)&=0,\label{eq:Jc}\\
 D_k(\tau)&=\frac{B_{2k}\tau^{2k-1}}{(2k)!}
 \left(\Li_{2-2k}(\e^{-2\tau})-2\Li_{2-2k}(\e^{-\tau})\right)
 -\boldsymbol1_{k=1}\frac\tau{24}.
 \label{eq:Dk}
\end{align}
The factor $-\tau/24$ must be kept in $D_1$; it comes from eta modularity,
not from the finite-size kernel.

\begin{theorem}[All orders uniformly through $\tau=0$]
\label{thm:uniform}
The functions in \eqref{eq:Ac}--\eqref{eq:Dk} extend real-analytically to
$\tau=0$. For $k\ge1$,
\begin{equation}
 D_k(0)=\frac{B_{2k}}{2k(2k-1)}(2^{1-2k}-2).
 \label{eq:Dzero}
\end{equation}
For every fixed $T>0$ and integer $K\ge0$, uniformly for $0\le\tau\le T$
as $x\to\infty$,
\begin{align}
 F_{\tau/x}(x)={}&x\Acal(\tau)-\frac12\log(\pi x)+\Jcal(\tau)
 +\sum_{k=1}^K\frac{D_k(\tau)}{x^{2k-1}}
 +\Ecal(\tau/x)+O_{T,K}(x^{-2K-1}).
 \label{eq:uniform}
\end{align}
For $K=0$ the term $-\tau/(24x)$ may be absorbed in the stated remainder.
For $K\ge1$ the remainder is exactly $(-1)^{K+1}r_K(\tau/x,\tau)$.
The endpoint interpretation is $F_0(x)=\log H_0(x)$ and $\Ecal(0)=0$.
\end{theorem}
\begin{proof}
The derivative \eqref{eq:Sder} is analytic at zero, and $S(0)=0$;
therefore $S(\tau)/\tau$ has the asserted removable singularity.
Also $(\tau/2)\coth(\tau/2)$ is a positive even analytic function with
value one at zero, so \eqref{eq:Jc} is analytic there.

For an integer $p\ge0$,
\[
 \Li_{-p}(\e^{-\tau})=
 \left(-\frac{\mathrm d}{\mathrm d\tau}\right)^p\frac1{\e^\tau-1}.
\]
The Laurent expansion at zero has leading term $p!\tau^{-p-1}$;
there are no other negative powers. Taking $p=2k-2$ in \eqref{eq:Dk}
proves removability and \eqref{eq:Dzero}. The elementary simplification
\[
 D_1(\tau)=\frac\tau{24}\left(\coth\tau-2\coth\frac\tau2\right)
\]
also displays the cancellation of its apparent linear term.

For $\tau>0$, rearranging \eqref{eq:forwardK} with $h=\tau/x$ gives
\eqref{eq:uniform}. Define
\[
 b_{K+1}(\tau)=\tau^{2K+1}a_{K+1}(\tau).
\]
The same Laurent argument shows that $b_{K+1}$ has a finite continuous
limit at zero. Thus $b_{K+1}$ is bounded on $[0,T]$, and
\eqref{eq:nextterm} supplies the required uniform constant. For $K=0$,
include the bounded term $\tau/(24x)$ separately. For fixed $x$, pass to
$\tau=0$ using \eqref{eq:qgamma}; alternatively, the Riemann-sum limit
below passes directly in the exact kernel identity.
\end{proof}

For example,
\begin{align*}
 \Acal(\tau)&=2\log2+\frac\tau2+\frac{\tau^2}{12}
                          -\frac{\tau^4}{480}+O(\tau^6),\\
 D_1(0)&=-\frac18, & D_2(0)&=\frac1{192}, &D_3(0)&=-\frac1{640},\\
 D_4(0)&=\frac{17}{14336},&D_5(0)&=-\frac{31}{18432},
 &D_6(0)&=\frac{691}{180224}.
\end{align*}
Hence the expansion reduces to the logarithmic central-binomial Stirling
series at the endpoint, with the correct normalization
$2x\log2-\tfrac12\log(\pi x)$.

\begin{proposition}[Continuous exact endpoint remainder]
\label[proposition]{prop:endrem}
For $x>0$ and $p=2K$, set $L=2\pi x$. The continuous extension of
$r_K(\tau/x,\tau)$ at $\tau=0$ is
\begin{equation}
 r_K^{(0)}(x)=\frac1{2\pi^2xL^p}
 \sum_{n\ge1}\frac1{n^{p+2}}
 \int_0^\infty
 \frac{(2\e^{-u}-\e^{-2u})u^p}{1+(u/(nL))^2}\dd u.
 \label{eq:endrem}
\end{equation}
It is strictly positive, satisfies the same first-omitted-term bound, and
is strictly log-convex as a sequence in $K$.
\end{proposition}
\begin{proof}
After the change $u=m\tau$, \eqref{eq:rK} is the corresponding Riemann sum
with mesh $\tau$. The exponential decay gives an integrable dominating
bound for each fixed $K,x$; the $n$-sum is dominated by $\sum n^{-p-2}$.
This proves the limit. The inequalities pass to the integral (strictness
also follows inside it). Rewriting the integrand as a positive measure
against $(u/(nL))^{2K}$ proves strict log-convexity exactly as in
\cref{lem:logconvex}.
\end{proof}

\begin{remark}[Scope of the endpoint assertion]
The uniform estimate is a large-$x$ statement for the whole closed range
$0\le hx\le T$. It does not require $hx$ to remain bounded away from zero,
and it is not merely a fixed-$x$ continuity statement for $q$-gamma.
Conversely, it does not assert uniformity as $T\to\infty$; the proof and
optimal-truncation constants below allow dependence on $T$.
\end{remark}

\section{Borel geometry and the cancellation of even axial poles}
\label{sec:borel}

\subsection{Convention and an explicit meromorphic transform}
For a formal series without constant term our convention is
\begin{equation}
 \Borel\!\left[\sum_{j\ge0}c_jh^{j+1}\right](\xi)
 =\sum_{j\ge0}\frac{c_j\xi^j}{j!},\qquad
 \mathcal L_0 b(h)=\int_0^\infty\e^{-\xi/h}b(\xi)\dd\xi.
 \label{eq:borelconv}
\end{equation}
Thus there is no extra $1/h$ in the Laplace integral. The perturbative
correction in \eqref{eq:forwardK}, apart from phase and amplitude, is
\begin{equation}
 \widetilde P(h,\tau)=-\frac h{24}
             +\sum_{k\ge1}(-1)^ka_k(\tau)h^{2k-1}.
 \label{eq:Pformal}
\end{equation}
The word perturbative here includes its full divergent Bernoulli series;
it excludes the separately prescribed modular term $\Ecal(h)$.

\begin{theorem}[Borel transform, pole divisor, and exact sum]
\label{thm:borel}
For fixed $\tau>0$, define
\begin{equation}
 \beta(\xi;\tau)=\frac1{4\pi^2}\sum_{n\ge1}\frac1{n^2}
 \left\{\frac1{\e^{\tau+\ii\xi/(2\pi n)}-1}
             +\frac1{\e^{\tau-\ii\xi/(2\pi n)}-1}\right\}.
 \label{eq:beta}
\end{equation}
Then the Borel transform of \eqref{eq:Pformal} is the meromorphic function
\begin{align}
 B_\tau(\xi)&=-\frac1{24}+\beta(\xi;2\tau)-2\beta(\xi;\tau)\label{eq:Btau}\\
 &=-\frac1{24}-\frac1{2\pi^2}\sum_{n\ge1}\frac1{n^2}
 \left\{\frac1{\e^{\tau+\ii\xi/(2\pi n)}-\epsilon_n}
             +\frac1{\e^{\tau-\ii\xi/(2\pi n)}-\epsilon_n}\right\},\notag\\
 \epsilon_n&=\begin{cases}1,&n\text{ odd},\\-1,&n\text{ even}.
                    \end{cases}\notag
\end{align}
All poles are simple. With $\ell\in\mathbb Z$, they are
\begin{equation}
 \begin{cases}
 4\pi^2 n\ell\ \pm 2\pi\ii n\tau,&n\text{ odd},\\
 2\pi^2 n(2\ell+1)\ \pm 2\pi\ii n\tau,&n\text{ even}.
 \end{cases}
 \label{eq:poles}
\end{equation}
At an upper-half-plane pole the residue is $\ii/(\pi n)$ for odd $n$
and $-\ii/(\pi n)$ for even $n$; the lower residues have the opposite
signs. In particular, the nearest poles are $\pm2\pi\ii\tau$ and every
even-indexed pure imaginary candidate has cancelled.

The series in \eqref{eq:Btau} converges normally away from its poles.
There are no poles on the positive real axis, and for all $h>0$,
\begin{equation}
 \mathcal L_0 B_\tau(h)=-\frac h{24}-R(h,\tau).
 \label{eq:exactborelsum}
\end{equation}
Consequently the phase, amplitude, and positive-direction Borel sum of the
Bernoulli block reconstruct $F_h(x)-\Ecal(h)$ exactly.
\end{theorem}
\begin{proof}
Expand the exponentials in \eqref{eq:beta} geometrically for real $\xi$.
Their sum is
\[2\sum_{m\ge1}\e^{-m\tau}\cos(m\xi/(2\pi n)).\]
The identity
\[
 \int_0^\infty\e^{-\xi/h}\cos\!\left(\frac{m\xi}{2\pi n}\right)\dd\xi
 =\frac{h}{1+(mh/(2\pi n))^2}
\]
shows that its Laplace transform is the last double sum in
\eqref{eq:Akernel}. Interchanging sums and integral is legitimate because
$\sum_n n^{-2}\sum_m\e^{-m\tau}<\infty$ and the Laplace weight is
integrable. This proves \eqref{eq:exactborelsum}.

The Taylor expansion of the cosine gives the Borel transform of the
Bernoulli coefficients. Equivalently, comparison with
\eqref{eq:coth} identifies those coefficients directly. The nearest
poles give a positive Taylor radius, so these coefficient identifications
are analytic identities near zero, not merely manipulations of a divergent
Taylor series in $h$.

To obtain the parity form, reindex the terms of $\beta(\xi;2\tau)$ by
$n=2j$ and use
\[
 \frac1{\e^{2u}-1}
 =\frac12\left(\frac1{\e^u-1}-\frac1{\e^u+1}\right).
\]
Its even-index Bose denominator cancels against $-2\beta(\xi;\tau)$,
leaving the stated alternating denominator. Solving for its zeros gives
\eqref{eq:poles}. At an upper odd pole the derivative of the denominator
is $\ii/(2\pi n)$, and at an upper even pole it is
$-\ii/(2\pi n)$; multiplication by $-1/(2\pi^2n^2)$ gives the residues.
The two signs of the imaginary part do not overlap, and different positive
$n$ have different imaginary parts. Thus no further pole collisions occur.

On a compact set in the $\xi$-plane, the large-$n$ denominators tend
uniformly to $\e^\tau-\epsilon_n$, both bounded away from zero. The factor
$n^{-2}$ proves normal convergence. This also proves that the displayed
pole list is complete.
\end{proof}

\subsection{A precisely oriented residue jump}
Only odd $n$ contribute on the positive imaginary axis. Let two admissible
Laplace contours pass on opposite sides of this ray. Fix their oriented
difference to enclose its poles by \emph{positive} small circles, so that
the residue theorem supplies $+2\pi\ii$ times each residue. In a common
convergence sector put
\[
 \zeta=\exp(-2\pi\ii\tau/h),\qquad |\zeta|<1.
\]
Then the perturbative contour difference is
\begin{equation}
 \Delta P=2\pi\ii\sum_{\substack{n\ge1\\n\text{ odd}}}
 \frac{\ii}{\pi n}\zeta^n
 =-2\sum_{n\text{ odd}}\frac{\zeta^n}{n}
 =\log\frac{1-\zeta}{1+\zeta}.
 \label{eq:jump}
\end{equation}
The logarithm is the branch analytic at $\zeta=0$ with value zero there.
One can first deform contours past finitely many poles and then pass to
the absolutely convergent residue sum; reversing the orientation reverses
the sign. The ray has a gap from the other rays in \eqref{eq:poles}, so
no undisclosed poles enter this deformation.

Equation \eqref{eq:jump} concerns the Borel sum of
\eqref{eq:Pformal}. It is not a claim that the original product and its
eta factor possess an unrestricted common analytic continuation across
$\Re h=0$. Keeping this boundary avoids identifying a valid local
residue calculation with a substantially stronger global continuation
theorem.

\subsection{The classical Borel transform emerges uniformly}
Use the actual endpoint expansion variable $\varepsilon=1/x$, so
$h=\tau\varepsilon$. Its Borel transform is
\begin{equation}
 \mathscr B_\tau(s)=\tau B_\tau(\tau s).
 \label{eq:scaledB}
\end{equation}

\begin{proposition}[Scaled Borel endpoint]
\label[proposition]{prop:scaledborel}
As $\tau\downarrow0$, locally uniformly away from the indicated poles,
\begin{equation}
 \mathscr B_\tau(s)\longrightarrow
 \mathscr B_0(s)=-4\sum_{\substack{n\ge1\\n\text{ odd}}}
                       \frac1{s^2+(2\pi n)^2}.
 \label{eq:Bzero}
\end{equation}
For every $T>0$ and $s\ge0$,
\begin{equation}
 |\mathscr B_\tau(s)|\le \frac T{24}+\frac16,
 \qquad 0<\tau\le T.
 \label{eq:Bbound}
\end{equation}
In particular, the positive-direction Laplace sums pass to the classical
endpoint by dominated convergence. The nearest scaled poles are always
$\pm2\pi\ii$, explaining the scale that appears in optimal truncation.
\end{proposition}
\begin{proof}
For odd $n$, the two fractions multiplied by $\tau$ tend to
$(1+\ii s/(2\pi n))^{-1}+(1-\ii s/(2\pi n))^{-1}$.
For even $n$, both tend to zero. The odd expression, multiplied by
$-1/(2\pi^2n^2)$, is exactly the summand in \eqref{eq:Bzero}.
Uniform domination by $C n^{-2}$ on compact sets away from the limit
poles proves local normal convergence. Poles with nonzero real part in
\eqref{eq:poles}, after division by $\tau$, escape to infinity.

For real $s$,
$|\e^{\tau\pm\ii\tau s/(2\pi n)}-\epsilon_n|\ge\e^\tau-1$.
Since $\tau/(\e^\tau-1)\le1$, summation yields
$\tau/24+\zeta(2)/\pi^2\le T/24+1/6$. The Laplace weight
$\e^{-xs}$ is integrable, which proves the final convergence statement.
\end{proof}

\begin{remark}[The missing flat term is not a summation ambiguity]
The positive real Borel ray has no singularity and its sum is unambiguous.
Nevertheless the exact normalized answer is that sum plus $\Ecal(h)$.
The formal power series alone cannot determine an independently added
flat function. Here the product normalization and eta identity determine
it. This example is therefore a useful concrete distinction between
uniqueness of a chosen Borel sum and uniqueness among all functions sharing
the same Poincar\'e expansion.
\end{remark}

\section{Sharp optimal truncation, including the leading constant}
\label{sec:optimal}
The nearest Borel singularity suggests the exponential size, but does not
by itself prove the actual minimum or its prefactor. We prove both from
the positive remainder.

\subsection{A uniform discrete saddle estimate}
We use the notation
\begin{equation}
 p=2K,\qquad L=2\pi x,\qquad
 g_{n,L}(u)=\frac1{1+(u/(nL))^2}.
 \label{eq:pL}
\end{equation}
A direct rescaling of \eqref{eq:rK} gives
\begin{equation}
 r_K(\tau/x,\tau)=\frac1{2\pi^2xL^p}
 \sum_{n\ge1}n^{-p-2}\,
 \tau\sum_{m\ge1}
 (2\e^{-m\tau}-\e^{-2m\tau})(m\tau)^p g_{n,L}(m\tau).
 \label{eq:rescaledr}
\end{equation}
At $\tau=0$ the inner Riemann sum means its integral as in
\eqref{eq:endrem}.

\begin{lemma}[Uniform gamma-saddle approximation]
\label[lemma]{lem:saddle}
Fix $T>0$ and $0<c<C<\infty$. As $p\to\infty$, uniformly for
$0\le\tau\le T$ and $c\le p/L\le C$,
\begin{equation}
 r_K(\tau/x,\tau)=
 \frac{p!}{\pi^2xL^p\bigl(1+(p/L)^2\bigr)}
 \left(1+O_{T,c,C}(p^{-1})\right).
 \label{eq:saddle}
\end{equation}
\end{lemma}
\begin{proof}
We first justify a relative, not merely absolute, Riemann-sum estimate.
For a twice differentiable, decaying function whose endpoint derivatives
and endpoint values vanish, the first periodic-Bernoulli remainder in the
trapezoidal formula gives
\begin{equation}
 \left|\tau\sum_{m\ge1}f(m\tau)-\int_0^\infty f(u)\dd u\right|
 \le\frac{\tau^2}{12}\int_0^\infty|f''(u)|\dd u.
 \label{eq:riemannbound}
\end{equation}
For completeness this follows by applying the finite-interval
Euler--Maclaurin identity through the $B_2$ term, letting its upper endpoint
go to infinity, and using $\sup_{0\le t\le1}|B_2(t)|=1/6$.
Here $f(u)=f_0(u)g_{1,L}(u)$ and $f_0(u)=\e^{-u}u^p$, with $p>2$,
so all the required endpoint terms vanish.

Normalize $f_0$ by $p!$ and let $U$ have this gamma density. Since
$f_0'/f_0=p/u-1$, the exact inverse moments of $U$ give
\begin{align}
 \int_0^\infty|f_0'(u)|\dd u&\le\frac{p!}{\sqrt{p-1}},\label{eq:fp}\\
 \int_0^\infty|f_0''(u)|\dd u&\le\frac{2p!}{p-1}.
 \label{eq:fpp}
\end{align}
Indeed, $\mathbb E(p/U-1)^2=1/(p-1)$ and
$\mathbb E(p/U^2)=1/(p-1)$; Cauchy--Schwarz and the formula for
$f_0''$ prove the bounds. The derivatives of $g_{1,L}$ satisfy
$\|g'_{1,L}\|_\infty=O(L^{-1})$ and
$\|g''_{1,L}\|_\infty=O(L^{-2})$. The product rule, with $L\asymp p$,
therefore gives
\begin{equation}
 \int|f''|\le C_1\frac{p!}{p}.
 \label{eq:fproduct}
\end{equation}
The integral $\int f=p!\,\mathbb E g_{1,L}(U)$ is bounded below by a
positive constant times $p!$ when $p/L\in[c,C]$: for example, the event
$U\le2(p+1)$ has probability at least $1/2$ by Markov's inequality,
and $g_{1,L}$ has a uniform positive lower bound on that event.
Equations \eqref{eq:riemannbound} and \eqref{eq:fproduct} prove a relative
$O_T(p^{-1})$ error, also uniformly as $\tau\downarrow0$.

The $n\ge2$ contribution in \eqref{eq:rescaledr} is
$O(2^{-p})$ relative to the $n=1$ contribution: discard its kernel by
$g_{n,L}\le1$, use the just-established upper moment bound, and sum
$\sum_{n\ge2}n^{-p-2}=O(2^{-p})$. The negative exponential
$\e^{-2u}$ contributes only $O_T(2^{-p})$ relatively. To see this without
assuming a fixed mesh, put $v=2u$ and apply the same Riemann estimate at
mesh $2\tau\le2T$; its unweighted integral is $p!/2^{p+1}$.
Hence \eqref{eq:rescaledr} reduces to
\begin{equation}
 \frac1{\pi^2xL^p}\int_0^\infty
       \frac{\e^{-u}u^p}{1+(u/L)^2}\dd u
       \left(1+O_T(p^{-1})\right).
 \label{eq:reducedsaddle}
\end{equation}
Finally, $\mathbb EU=p+1$ and $\operatorname{Var}U=p+1$.
Taylor's theorem applied to $g_{1,L}$ about $p$ gives
\[
 \mathbb E g_{1,L}(U)
 =g_{1,L}(p)+O(L^{-1})+O((p+2)L^{-2})
 =\frac1{1+(p/L)^2}+O(p^{-1}).
\]
The leading value is bounded away from zero on $[c,C]$, so this is a
relative estimate. Substitution into \eqref{eq:reducedsaddle} proves
\eqref{eq:saddle}. The integral case $\tau=0$ is included throughout.
\end{proof}

\subsection{Locating the true minimum}
\begin{theorem}[Uniform sharp minimum]
\label{thm:optimal}
For every $T>0$, uniformly for $0\le\tau\le T$ as $x\to\infty$,
\begin{equation}
 \boxed{\displaystyle
 \min_{K\ge0}r_K(\tau/x,\tau)
 =\frac{\e^{-2\pi x}}{\pi\sqrt{x}}
       \left(1+O_T(x^{-1})\right).}
 \label{eq:optimal}
\end{equation}
Every minimizing index satisfies
\begin{equation}
 K=\pi x+O_T(1).
 \label{eq:optimalK}
\end{equation}
More generally, the asymptotic expression in \eqref{eq:optimal} holds
for any integer $K=\pi x+O(1)$, uniformly when the implied displacement
bound is fixed.
\end{theorem}
\begin{proof}
Take the quotient of \eqref{eq:saddle} at $p+2$ and $p$. Uniformly when
$p/L$ stays in a fixed compact subinterval of $(0,\infty)$,
\begin{align}
 \frac{r_{K+1}}{r_K}
 &=\frac{(p+1)(p+2)}{L^2}
   \frac{1+(p/L)^2}{1+((p+2)/L)^2}
   \left(1+O_T(p^{-1})\right)\notag\\
 &=\left(\frac pL\right)^2\left(1+O_T(L^{-1})\right)
 \quad\text{when }p\asymp L.
 \label{eq:quotient}
\end{align}
Choose a sufficiently large fixed $C_T$. At an even integer $p$ within
two of $L-C_T$, the quotient is less than one; at an even integer within
two of $L+C_T$, it is greater than one. This follows by expanding
$(1\pm C_T/L)^2$, with $C_T$ larger than the uniform constant in the error
term. Strict log-convexity from \cref{lem:logconvex,prop:endrem} then puts
every global minimum between these two indices. This is the step that
upgrades a local least-term heuristic to a statement about all $K\ge0$.

For $p=L+O_T(1)$, Stirling's formula gives
\begin{equation}
 \frac{p!}{L^p}=\sqrt{2\pi L}\,\e^{-L}
                 \left(1+O_T(L^{-1})\right).
 \label{eq:stirlingsaddle}
\end{equation}
Also $1+(p/L)^2=2+O_T(L^{-1})$. Since
$\sqrt{2\pi L}/(2\pi^2x)=1/(\pi\sqrt{x})$, substitution into
\eqref{eq:saddle} proves \eqref{eq:optimal}. The same computation proves
the assertion for any bounded displacement from $\pi x$.
\end{proof}

\begin{remark}[What fixes the universal constant]
Three effects contribute. The $n=1$ partial-fraction kernel dominates;
$2\e^{-u}$ dominates $2\e^{-u}-\e^{-2u}$ at the large-order saddle; and
the kernel denominator is asymptotically $1+(p/L)^2=2$. Its factor $1/2$
is essential. Replacing the exact remainder by the first omitted term
would give a different prefactor and would not prove the sharp minimum.
The finite mesh becomes invisible to relative order $O_T(1/x)$ because
the gamma saddle has width of order $\sqrt{x}$ while $\tau\le T$.
\end{remark}

\section{A sharp modular-visibility boundary}
\label{sec:visibility}

\begin{lemma}[Elementary bounds for the modular series]
\label[lemma]{lem:modbounds}
For $0<z<1$, writing $\Ecal(z)=\sum_{j\ge1}\sigma_{-1}(j)z^j$,
\begin{align}
 z\le\Ecal(z)&\le\frac{z}{(1-z)^2},\label{eq:modbound}\\
 0\le\Ecal(z)-\sum_{j=1}^M\sigma_{-1}(j)z^j
 &\le\frac{z^{M+1}((M+1)-Mz)}{(1-z)^2}.
 \label{eq:modtail}
\end{align}
In particular $\Ecal(z)=z+O(z^2)$.
\end{lemma}
\begin{proof}
The coefficient is positive, equals one at $j=1$, and satisfies
$\sigma_{-1}(j)\le j$. Sum $\sum jz^j$ and its tail.
\end{proof}
The notation $\Ecal(z)$ in this lemma is its generating function; elsewhere
$\Ecal(h)$ means its value at $z=\e^{-4\pi^2/h}$.

\begin{theorem}[Visibility threshold and resolved boundary layer]
\label{thm:visibility}
Let $r_*(x,\tau)=\min_K r_K(\tau/x,\tau)$. For fixed $\tau>0$,
\begin{equation}
 \frac{r_*(x,\tau)}{\Ecal(\tau/x)}
 =\frac1{\pi\sqrt{x}}
  \exp\!\left[x\left(\frac{4\pi^2}{\tau}-2\pi\right)\right]
  \left(1+O_\tau(x^{-1})\right).
 \label{eq:ratio}
\end{equation}
Thus the ratio tends to infinity for $0<\tau<2\pi$, to zero for
$\tau>2\pi$, and is asymptotic to $1/(\pi\sqrt{x})$ at $\tau=2\pi$.
For fixed real $s$, define
\begin{equation}
 \tau_x(s)=2\pi-\frac{\log x}{2x}+\frac{s}{x}.
 \label{eq:boundary}
\end{equation}
Then, uniformly for $s$ in compact subsets of $\R$,
\begin{equation}
 \boxed{\displaystyle
 \frac{r_*(x,\tau_x(s))}{\Ecal(\tau_x(s)/x)}
 \longrightarrow \frac{\e^{-s}}\pi.}
 \label{eq:boundarylimit}
\end{equation}
The same limit holds with any $K=\pi x+O(1)$ in place of the minimizer.
\end{theorem}
\begin{proof}
Combine \eqref{eq:optimal} with \eqref{eq:modbound}. In the boundary layer
put $\delta=\tau_x-2\pi$. The exact expansion of the reciprocal gives
\[
 x\left(\frac{4\pi^2}{2\pi+\delta}-2\pi\right)
 =-x\delta+O(x\delta^2)
 =\frac12\log x-s+O\!\left(\frac{(\log x)^2}{x}\right).
\]
The error is uniform for bounded $s$. Substitution proves
\eqref{eq:boundarylimit}; \cref{thm:optimal} supplies the last assertion.
\end{proof}

The equal-size boundary occurs, to this accuracy, at
\begin{equation}
 \tau=2\pi-\frac{\log x}{2x}-\frac{\log\pi}{x}+o(x^{-1}).
 \label{eq:equalboundary}
\end{equation}
More generally, to make the best truncation residual a fraction $\eta>0$
of the leading modular signal in the boundary scaling, the limiting
condition is $s>-\log(\pi\eta)$. This is a quantitative precision budget,
not merely a comparison of two exponentials with their polynomial
prefactors discarded.

\subsection{Interpretation and limitations}
Suppose one approximates the nonmodular carrier by truncating its Bernoulli
series, retaining phase and amplitude exactly. The smallest possible
absolute forward error among such approximants is $r_*$. For
$\tau<2\pi$ fixed, every such truncation has error much larger than the
modular correction. Simply appending a modular term to an algebraically
truncated expression therefore does not resolve that term in absolute
accuracy. At and above the threshold, near-optimal truncation does resolve
it asymptotically.

This is not a positive-axis Stokes phenomenon: \cref{thm:borel} shows that
there are no Borel poles on that axis. Nor is it a claim that a real
exponential with action $2\pi$ must be inserted as an independent sector
in the exact answer. It is the scale of the error left by optimal
truncation of a Borel-summable block whose nearest poles are imaginary.
Exact integration of that block, or direct evaluation of the positive
remainder, can achieve accuracy below $r_*$. The theorem is sharp for its
stated family of approximants, not for every conceivable algorithm.

\section{Certified inversion and a convergent modular inverse series}
\label{sec:inverse}
All derivatives in this section are taken with $h$ \emph{fixed}. Although
we use $\tau=hx$ to state uniform bounds, differentiating along the curve
$h=\tau/x$ with $\tau$ fixed would be a different operation and would give
the wrong inverse sensitivity.

\subsection{A uniform lower bound for the exact slope}
\begin{theorem}[Monotonicity and quantitative inverse sensitivity]
\label{thm:slope}
For $h>0$ and $x>0$, the derivative is
\begin{equation}
 F'_h(x)=2hx+2h\sum_{m\ge1}
           \frac{\e^{-mhx}-\e^{-2mhx}}{\e^{hm}-1}>0.
 \label{eq:slope}
\end{equation}
Writing $\tau=hx$,
\begin{equation}
 S'(\tau)-\frac{h}{2\sinh\tau}
 \le F'_h(x)\le S'(\tau),
 \qquad
 F'_h(x)\ge S'(\tau)-\frac1{2x}.
 \label{eq:slopebounds}
\end{equation}
For $x\ge1$ this gives the uniform positive constant
\begin{equation}
 F'_h(x)\ge d_0:=2\log2-\frac12>0.
 \label{eq:d0}
\end{equation}
The bounds extend to $h=0$ by continuity. If $F_h(x)=L$ and
$x,\widehat x\ge1$, then
\begin{equation}
 |x-\widehat x|\le\frac{|F_h(\widehat x)-L|}{d_0}.
 \label{eq:certificate}
\end{equation}
\end{theorem}
\begin{proof}
Differentiate the absolutely locally convergent series
\eqref{eq:exactforward}. Each summand in \eqref{eq:slope} is positive.
For $v>0$, \eqref{eq:coth} and $\e^v\ge1+v$ imply
\[
 \frac1v-\frac12\le\frac1{\e^v-1}\le\frac1v.
\]
The upper comparison in \eqref{eq:slope} sums to
$2\tau+2\log(1+\e^{-\tau})=S'(\tau)$. The loss in the lower comparison is
\[
 h\sum_{m\ge1}(\e^{-m\tau}-\e^{-2m\tau})
 =\frac{h}{2\sinh\tau}\le\frac1{2x}.
\]
This proves \eqref{eq:slopebounds} and \eqref{eq:d0}. For fixed $x$, the sum in \eqref{eq:slope} is a Riemann sum
with mesh $h$ for the continuous, integrable function
$(\e^{-xt}-\e^{-2xt})/(\e^t-1)$. Its limit is
$2\int_0^\infty(\e^{-xt}-\e^{-2xt})/(\e^t-1)\dd t
=2\psi(2x+1)-2\psi(x+1)=F_0'(x)$,
which proves the endpoint derivative bound without differentiating a
merely pointwise limit. Finally the mean value
theorem on the interval between $x$ and $\widehat x$ proves
\eqref{eq:certificate}.
\end{proof}

A certificate should also certify existence, not merely an error
conditional on a root. Suppose a computable approximation $T_h(t)$ has
known bounds $|F_h(t)-T_h(t)|\le\epsilon(t)$ at two points $1\le a<b$.
If
\begin{equation}
 T_h(a)+\epsilon(a)\le L\le T_h(b)-\epsilon(b),
 \label{eq:existbracket}
\end{equation}
continuity and strict monotonicity give one and only one root in $[a,b]$.
At an interior trial point, \eqref{eq:certificate} then gives
\[
 |x-\widehat x|\le
 \frac{|T_h(\widehat x)-L|+\epsilon(\widehat x)}{d_0}.
\]
The positive kernel and \eqref{eq:modtail} supply such $\epsilon$ without
making an asymptotic remainder into an unsupported numerical assertion.

\subsection{Inverting the exact nonmodular carrier}
Define
\begin{equation}
 F_h^0(x)=F_h(x)-\Ecal(h),\qquad G_h^0=(F_h^0)^{-1}.
 \label{eq:carrier}
\end{equation}
By \cref{thm:borel}, $F_h^0$ consists of the exact phase and amplitude and
the Borel-summed Bernoulli block. It is not a finite truncation of that
block. Since $\Ecal(h)$ is independent of $x$, both carriers have the
same derivative. If $x_0=G_h^0(L)$, then the exact inverse relation is
\begin{equation}
 F_h^{-1}(L)=G_h^0(L-\Ecal(h)).
 \label{eq:inverseexact}
\end{equation}
A shift in target therefore encodes every modular inverse sector.

\begin{theorem}[Uniformly convergent modular inverse expansion]
\label{thm:inverse}
Fix $T>0$. Let $x_0\ge1$, $h>0$, $hx_0\le T$, and $L=F_h^0(x_0)$.
Put
\begin{equation}
 M_T=2T+8,\quad
 r=\min\left\{\frac14,\frac{d_0}{4M_T}\right\},\quad
 \rho=\frac{d_0r}{2},\qquad
 \Dcal_h=\frac1{F'_h(x)}\frac{\mathrm d}{\mathrm dx}.
 \label{eq:radii}
\end{equation}
For $\Ecal(h)<\rho$,
\begin{equation}
 F_h^{-1}(L)=x_0+
 \sum_{n\ge1}\frac{(-1)^n\Ecal(h)^n}{n!}
 \left[\Dcal_h^{\,n-1}\frac1{F'_h(x)}\right]_{x=x_0}.
 \label{eq:inverseE}
\end{equation}
The series converges absolutely, and its tail after $n=N$ is at most
\begin{equation}
 r\,\frac{(\Ecal(h)/\rho)^{N+1}}{1-\Ecal(h)/\rho}.
 \label{eq:inverseEtail}
\end{equation}
It may also be collected into a convergent series in $z=\e^{-4\pi^2/h}$:
\begin{equation}
 F_h^{-1}(L)=x_0+\sum_{m\ge1}C_m(h,x_0)z^m,
 \label{eq:inversez}
\end{equation}
where every coefficient is the finite expression
\begin{align}
 C_m(h,x_0)={}&\sum_{n=1}^m\frac{(-1)^n}{n!}
 \left[\Dcal_h^{\,n-1}\frac1{F'_h(x)}\right]_{x=x_0}\notag\\
 &\hspace{8mm}\times
 \sum_{\substack{j_1+\cdots+j_n=m\\j_i\ge1}}
        \prod_{i=1}^n\sigma_{-1}(j_i).
 \label{eq:Cm}
\end{align}
If $0<r_z<1$ satisfies $r_z/(1-r_z)^2<\rho$, then for $z<r_z$ the tail
of \eqref{eq:inversez} after degree $M$ is at most
\begin{equation}
 r\,\frac{(z/r_z)^{M+1}}{1-z/r_z}.
 \label{eq:inverseztail}
\end{equation}
\end{theorem}
\begin{proof}
The product logarithms, hence $F_h^0$, are holomorphic when $\Re x>0$;
the series provides a single-valued branch there. For $|w|\le1/4$,
$\Re(x_0+w)\ge3/4$. Differentiating \eqref{eq:slope} once and estimating the result gives
\begin{align*}
 |F_h''(x_0+w)|
 &\le 2h+2h^2\sum_{m\ge1}
     \frac{m(\e^{-mh a}+2\e^{-2mh a})}{\e^{hm}-1}\\
 &\le2h+2h\left(\frac1{\e^{ha}-1}
                    +\frac2{\e^{2ha}-1}\right)
 \le2h+\frac4a\le M_T,
\end{align*}
where $a=\Re(x_0+w)$ and $h\le T$ follows from $x_0\ge1$.
Consequently, on $|w|=r$,
\[
 |F_h^0(x_0+w)-L-F'_h(x_0)w|\le\tfrac12M_Tr^2.
\]
For $|s|<\rho$, the right side plus $|s|$ is strictly less than
$d_0r\le|F'_h(x_0)w|$. Rouch\'e's theorem gives a unique root of
$F_h^0(x_0+w)=L+s$ in $|w|<r$.
Moreover $|F'_h(x_0+w)-F'_h(x_0)|\le M_Tr\le d_0/4$,
so the derivative never vanishes there. The roots form one holomorphic
inverse $G_h^0(L+s)$ on $|s|<\rho$, with
$|G_h^0(L+s)-x_0|<r$.

Cauchy's coefficient bound is therefore
$|[s^n](G_h^0(L+s)-x_0)|\le r\rho^{-n}$.
Repeated differentiation of an inverse gives
\[
 (G_h^0)^{(n)}(L)
 =\left[\Dcal_h^{\,n-1}(1/F'_h)\right]_{x_0}\quad(n\ge1).
\]
Taylor expansion at $s=-\Ecal(h)$ proves \eqref{eq:inverseE}, and the
geometric sum gives \eqref{eq:inverseEtail}.

For the collection into $z$, temporarily regard $z$ as an independent
complex variable, holding $h,x_0$ fixed in the carrier. By
\cref{lem:modbounds}, $|\sum\sigma_{-1}(j)z^j|<\rho$ for $|z|\le r_z$.
Absolute convergence of \eqref{eq:inverseE} and positivity of the
coefficients of $\Ecal(z)$ justify collecting powers at every positive
$z$ with $\Ecal(z)<\rho$. On the indicated disk, the analytic composition
is bounded by $r$ relative to $x_0$, and its coefficients are exactly \eqref{eq:Cm} by the Cauchy product. A second
Cauchy estimate gives \eqref{eq:inverseztail}. Evaluating the auxiliary
variable at its physical value $z(h)$ completes the proof.
\end{proof}

The first two coefficients illustrate the separate sources of nonlinear
mixing:
\begin{align}
 C_1&=-\frac1{F'_h(x_0)},\label{eq:C1}\\
 C_2&=-\frac{3}{2F'_h(x_0)}
      -\frac{F_h''(x_0)}{2(F'_h(x_0))^3}.
 \label{eq:C2}
\end{align}
The $3/2$ is the divisor sum $\sigma_{-1}(2)$; the second term is inverse
curvature. Suppressing either changes the second modular sector.

Uniformly for $hx_0\le T$ as $x_0\to\infty$ and $h\downarrow0$,
\begin{equation}
 F_h^{-1}(L)-x_0
 =-\frac{\e^{-4\pi^2/h}}{S'(hx_0)}
 \left(1+O_T(x_0^{-1})+O_T(\e^{-4\pi^2/h})\right).
 \label{eq:inverseleading}
\end{equation}
The shift is negative because adding a positive $x$-independent term to an
increasing function lowers its inverse at a fixed target.

\subsection{What the visibility theorem says for inverse accuracy}
For a small forward residual $\delta$, the exact local inverse displacement
obeys
\begin{equation}
 G_h^0(L+\delta)-x_0
 =\frac{\delta}{F'_h(x_0)}+O_T(\delta^2).
 \label{eq:residualtransport}
\end{equation}
In particular, compare the transported near-optimal residual of magnitude
$r_K$ with the modular target shift $\Ecal(h)$, at the same center. Their
inverse displacement ratio is
\begin{equation}
 \frac{|G_h^0(L\pm r_K)-x_0|}
 {|G_h^0(L-\Ecal(h))-x_0|}
 =\frac{r_K}{\Ecal(h)}\left(1+O_T(r_K+\Ecal(h))\right).
 \label{eq:inverseratio}
\end{equation}
Thus \cref{thm:visibility} persists under this exact local transport.
This does not assert an optimal-truncation theorem for a separately
constructed formal inverse series. Such a theorem would need its own
large-order analysis. The convergent series in \cref{thm:inverse} instead
uses the \emph{exact} carrier and so has a different, explicitly controlled
source of error.

\section{Reproducible computations and an executable error budget}
\label{sec:checks}

\subsection{Cancellation-free evaluation of a tiny remainder}
Subtracting a truncated series from a function can destroy all significant
digits at optimal truncation. The positive representation avoids that
problem. With $p=2K$, $s=p+2$, and
\[
 c=\frac{2h^{p+1}}{(2\pi)^{p+2}},\qquad
 W_m=(2\e^{-m\tau}-\e^{-2m\tau})m^p,
\]
it reads
\begin{equation}
 r_K=c\sum_{m\ge1}W_m
        \sum_{n\ge1}\frac{n^{-s}}{1+(mh/(2\pi n))^2}.
 \label{eq:computesum}
\end{equation}
Let $R_{M,N}$ be the positive partial sum through $m=M$, $n=N$.
If
\[
 \theta_M=\e^{-\tau}\left(\frac{M+2}{M+1}\right)^p<1,
\]
then the exact omitted-tail estimate is
\begin{align}
 0\le r_K-R_{M,N}\le c\Bigg(&\zeta(s,N+1)\sum_{m=1}^M W_m\notag\\
 &+\zeta(s)\frac{2\e^{-(M+1)\tau}(M+1)^p}{1-\theta_M}\Bigg).
 \label{eq:computetail}
\end{align}
Here $\zeta(s,N+1)=\sum_{n>N}n^{-s}$ is the Hurwitz-zeta tail.
To prove this, discard each kernel by one. The missing $n$-tail on retained
$m$ gives the first term. For omitted $m$, use $W_m\le2\e^{-m\tau}m^p$;
its successive ratios decrease once the first omitted index is reached,
and are at most $\theta_M$. Summing the resulting geometric majorant
gives the second term. No subtraction of nearly equal large quantities
occurs.

The supplied \texttt{verify.py} implements \eqref{eq:computesum} and
\eqref{eq:computetail}. It uses 110 decimal digits. Tail inequalities are
proved in exact real arithmetic in this article; the program evaluates
them with ordinary high-precision rounding, \emph{not} outward-rounded
interval arithmetic. Accordingly the reported runs are reproducible
checks, not independently certified interval computations.

\subsection{Actual results of the supplied run}
The full script was run in preparing this manuscript. Exact symbolic
arithmetic checks the six endpoint coefficients displayed after
\cref{thm:uniform}. Eight exact identities also compare the endpoint Borel Taylor
coefficients with the odd-pole lattice. Ten rational residue-cancellation
checks are included. Independent finite products at integer indices check
\eqref{eq:exactforward}; the four absolute discrepancies were all below
$3\times10^{-110}$. This is an independent normalization test because it
does not evaluate both sides from the same infinite-product tail.

\begin{table}[htbp]
\centering
\caption{Positive remainders checked against direct high-precision
subtraction at nonoptimal orders.}
\label{tab:positive}
\begin{tabular}{@{}rrrrr@{}}
\toprule
$x$&$\tau$&$K$&$r_K$&$r_K/(a_{K+1}h^{2K+1})$\\
\midrule
5&1&4&$7.91517695222\cdot10^{-10}$&0.9191817471\\
7&2&6&$1.21204088308\cdot10^{-13}$&0.9160370429\\
8&0.5&6&$2.17819186412\cdot10^{-14}$&0.9340864521\\
\bottomrule
\end{tabular}
\end{table}
Each remainder in \cref{tab:positive} was positive, had the proved sign in
the forward identity, and lay below the first omitted term. The positive
partial-sum bounds enclosed the directly computed value at the stated
working precision.

\begin{table}[htbp]
\centering
\caption{Near-optimal truncation at $\tau=1$. The script uses
$K=\operatorname{round}(\pi x)$, not an unverified claim to have found the
exact minimizing integer.}
\label{tab:optimal}
\begin{tabular}{@{}rrrr@{}}
\toprule
$x$&$K$&$r_K$&$r_K\pi\sqrt{x}\,\e^{2\pi x}$\\
\midrule
5&16&$3.17973793023\cdot10^{-15}$&0.9835348221\\
10&31&$5.21956022965\cdot10^{-29}$&1.0053351399\\
20&63&$1.88567218808\cdot10^{-56}$&0.9958300329\\
40&126&$3.55472089880\cdot10^{-111}$&0.9979153044\\
\bottomrule
\end{tabular}
\end{table}
The mathematical relative upper bounds on omitted positive summands in
these four evaluations, as evaluated by the program, are respectively
below $4.2\cdot10^{-31}$, $6.1\cdot10^{-39}$, $1.3\cdot10^{-39}$, and
$1.7\cdot10^{-38}$. This illustrates why a value below $10^{-110}$ can
still be computed \emph{relatively} at 110 digits: the computation uses
positive small terms, not subtraction from a number of order one.
At $x=20$, tests with $\tau=0.1$, $2\pi$, and $10$ give normalized ratios
0.9958300329, 0.9958300329, and 0.9958300327, respectively.

\begin{table}[htbp]
\centering
\caption{Resolved boundary layer. The same near-optimal rule for $K$ is
used. Finite-$x$ values need not approach their limits monotonically.}
\label{tab:boundary}
\begin{tabular}{@{}rrr@{}}
\toprule
$x$&ratio at $s=0$&ratio at $s=1$\\
\midrule
10&0.3269592979&0.1177674042\\
20&0.3227622917&0.1168425200\\
40&0.3220074088&0.1171886859\\
80&0.3214572071&0.1174603693\\
\midrule
Limit&$1/\pi=0.3183098862$&$\e^{-1}/\pi=0.1170996630$\\
\bottomrule
\end{tabular}
\end{table}
The error in the exponent of \eqref{eq:boundarylimit} is of order
$(\log x)^2/x$, so the visible approach in \cref{tab:boundary} is slower
than the $O(1/x)$ accuracy of the pure optimal-remainder formula.

For an inverse check, take $h=1.2$, $x_0=5$, and the exact target
$L=F_h(5)-\Ecal(h)$. The independently computed inverse shift and its
first-order prediction are
\begin{align*}
 F_h^{-1}(L)-5&=-4.29546689467745006556\cdot10^{-16},\\
 -\Ecal(h)/F'_h(5)&=-4.29546689467744988132\cdot10^{-16}.
\end{align*}
Their ratio is $1.00000000000000004289\ldots$, with the correct negative
sign. Separate slope checks at $(h,x)=(0.01,1),(0.1,20),(0.25,4),(1,3)$
satisfy both sides of \eqref{eq:slopebounds}.

\subsection{Using the results rather than only inspecting coefficients}
For a prescribed target $L$, a usable workflow is to choose a center from
the repository's monotone dilogarithmic phase, bracket the exact root by
\eqref{eq:existbracket}, and then refine using the exact slope. A finite
Bernoulli expansion is adequate when its first-omitted-term bound fits the
error budget. Near exponential precision, evaluate the positive remainder
or the exact Borel kernel and retain as many modular sectors as
\eqref{eq:modtail} requires. Equation \eqref{eq:certificate} converts that
forward budget into an inverse budget. The modular inverse series is useful
when an exact carrier center is already available; it is not a replacement
for obtaining that center accurately.

To reproduce the supplied computations, run
\begin{verbatim}
python verify.py
\end{verbatim}
from the package directory. This writes \texttt{verification\_results.json}.
The optional \texttt{--quick} switch omits the larger cases. Python 3.10
or later, mpmath, and SymPy are required. The source was run with mpmath
1.3.0 and SymPy 1.14.0. The script raises an exception if an asserted test
fails and uses no network service or proprietary computer algebra system.

\section{Further research questions}
\label{sec:questions}
The following questions separate extensions plausibly accessible from this
kernel method from broader problems requiring additional ideas. They are
proposed research targets, not claims that the corresponding statements
are already proved.

\subsection{Higher optimal-remainder coefficients and the minimizing integer}
\textbf{Question.} Can one compute a complete expansion for
$r_K\pi\sqrt{x}\e^{2\pi x}$ in inverse powers of $x$, uniformly for bounded
$K-\pi x$ and $0\le\tau\le T$, and determine the exact minimizing integer
outside an explicitly described tie set?
The proof of \cref{lem:saddle} keeps only a first relative error. Higher
gamma moments and a sharper discrete-sum estimate should expose the
integer-rounding dependence. It is important to distinguish a smooth
large-order expansion from exponentially small mesh effects, particularly
when two adjacent truncation orders nearly tie.

\subsection{The range of a growing crossover parameter}
\textbf{Question.} What is the largest growth range $T=T(x)$ for which the
universal constant in \eqref{eq:optimal} remains valid?
The proof gives an initial diagnostic: the mesh error from
\eqref{eq:riemannbound} is $O(\tau^2/x)$ near the saddle. Thus
$\tau=o(\sqrt{x})$ is a natural first range to investigate, but a sharp
statement needs uniform constants and may be improved by exploiting
analyticity rather than a twice-differentiable estimate. When the mesh is
comparable to the saddle width, a genuinely discrete correction may remain.
This is not covered by the fixed-$T$ theorem.

\subsection{General Gaussian quotients and positivity criteria}
\textbf{Question.} For products of shifted $q$-gamma functions with balanced
slopes, which combinations give a positive kernel analogous to
$2\e^{-u}-\e^{-2u}$? What replaces the optimal error constant when the
kernel changes sign?
For positive kernels, moment log-convexity may survive and make a global
minimum accessible. With sign-changing kernels, cancellation can defeat
the present lower-bound argument even if the nearest Borel singularities
are known. A useful classification should state both the slope data and
the kernel-positivity condition, rather than infer one from the other.

\subsection{Arithmetic descriptions of pole cancellations}
\textbf{Question.} Can pole cancellation in finite products of
$q$-Pochhammer factors be expressed as an explicit finite arithmetic
transform of the slope and shift multiplicities?
The central case reduces Bose denominators to alternating Bose/Fermi
ones and removes even axial poles. Rational slope ratios suggest a finite
residue-class calculation, whereas irrational slopes may create
nonperiodic pole configurations. Establishing a criterion would turn a
case-by-case cancellation into a reusable transseries operation.

\subsection{Resurgence of the independently reverted crossover series}
\textbf{Question.} Determine the Borel singularities and Stokes data of the
formal inverse obtained by reverting the full polynomial--logarithmic
crossover series, including its $\log h$ coefficients.
The exact-carrier construction in \cref{thm:inverse} does not answer this.
Composition, moving singularities, and nonlinear sector mixing must be
controlled. A satisfactory result should also identify which singularities
are inherited from the forward kernel and which can arise from zeros of
the analytically continued derivative.

\subsection{A hyperasymptotic algorithm below the visibility threshold}
\textbf{Question.} Can the positive remainder be expanded again, using the
nearest pole pair, to produce a rigorously certified algorithm resolving
$\Ecal(h)$ for fixed $\tau<2\pi$ with fewer operations than direct tail
summation?
There is no impossibility obstruction: exact kernel integration already
contains the missing information. The question is computational complexity
and stable error accounting. A useful result would specify the number of
retained pole pairs and arithmetic precision as functions of $h,\tau$ and
the requested relative accuracy of the modular term.

\subsection{Root-of-unity and complex-parameter crossover}
\textbf{Question.} Which versions of the endpoint and sharp-remainder
arguments persist near nontrivial roots of unity, where several Borel rays
and residue classes interact?
The real positive kernel is lost in general. The comparison with the
$q$-Pochhammer literature \cite{garoufalidiskashaev,fantinirella} suggests
natural complex variants but not a proof of the positivity-based minimum.
A first tractable task is a contour-specific remainder bound away from
coalescing rays, before attempting a uniform transition across them.

\subsection{Machine-checked inequalities and interval implementation}
\textbf{Question.} Can the exact kernel identity, finite geometric remainder,
log-convexity, and residual certificate be formalized as a small reusable
Lean module, independently of a full transseries datatype?
The finite geometric identity and positive-sum estimates are promising
first targets. Tonelli interchange, the endpoint Riemann estimate, and the
uniform gamma-saddle constants should then be separate theorems. In
parallel, an outward-rounded implementation of \eqref{eq:computetail}
would turn the present executable checks into numerical certificates.
Neither deliverable is included as already completed here.

\subsection{Recovering a modular sector from additional functional data}
\textbf{Question.} Which minimal functional relations, normalization, and
growth conditions determine the flat modular correction uniquely from the
Borel-summed carrier?
The exact eta identity supplies such data for this example. Abstractly,
the same asymptotic power series admits additions of flat functions.
A useful uniqueness theorem should state the extra hypotheses explicitly;
formal coefficients alone cannot provide them.

\section{Conclusion}
The main mechanism is elementary but stronger than a fixed-order
asymptotic expansion. A positive partial-fraction remainder remains exact
when the truncation order grows with the large parameter. That permits a
uniform endpoint theorem, a global optimal-truncation argument through
moment log-convexity, and a sharp comparison with an independently
normalized modular sector. The resulting transition occurs at $hx=2\pi$
with a logarithmic displacement of order $(\log x)/x$, and is preserved by
exact local inverse transport.

The Borel calculation explains the scale through a stable nearest pole
pair, while the remainder calculation fixes the prefactor and the true
minimum. The inverse construction then separates a divergent but exactly
summable carrier from a genuinely convergent modular correction series.
These distinctions are precisely what is needed to extend the inspected
repository chapter without conflating a formal all-order formula,
an asymptotic error estimate, a Borel sum, and an exact normalized function.

\appendix
\section{Source and claim audit}
\label{app:audit}
\begin{longtable}{@{}p{0.27\textwidth}p{0.66\textwidth}@{}}
\toprule
Item & Status in this manuscript\\
\midrule
\endfirsthead
\toprule
Item & Status in this manuscript\\
\midrule
\endhead
Exact interpolation and phase & Taken from the inspected ProveIt companion;
normalization and derivative rederived in \cref{sec:normalization}.\\[5pt]
Eta modular factor & Classical identity, already used in the companion;
not a novelty claim.\\[5pt]
Fixed-order Bernoulli coefficients & Already in the companion;
regularized at the endpoint here.\\[5pt]
Uniform $\tau=0$ matching & Proved in \cref{thm:uniform}, addressing the
explicit boundary of the companion's fixed-$\tau_0$ result.\\[5pt]
General $q$-product resurgence & Established antecedents
\cite{costingaroufalidis,garoufalidiskashaev,fantinirella}; not claimed new.\\[5pt]
Central Borel pole cancellations & Explicitly derived here, with signs,
residues, and scaled endpoint fixed. No literature-wide priority claim.\\[5pt]
Sharp minimum and visibility layer & Full proofs in
\cref{thm:optimal,thm:visibility}; the central quantitative extension
advanced for independent checking.\\[5pt]
Modular inverse sectors & Proved convergent about the exact carrier, with
explicit radii and tails. Does not prove general inverse resurgence.\\[5pt]
Numerical verification & Actual run of supplied code; exact symbolic
subchecks and 110-digit numerical checks. Not interval-certified.\\[5pt]
Lean status & No new Lean file supplied, and no claim of machine-checking
these analytic theorems.\\[5pt]
Repository coverage & Relevant companion ranges and READMEs inspected;
the full oversized canonical TeX was not read.\\
\bottomrule
\end{longtable}

\section{Reproduction and file inventory}
The package contains \texttt{article.tex}, the compiled \texttt{article.pdf},
\texttt{verify.py}, the actual \texttt{verification\_results.json}, a short
\texttt{README.md}, and pinned numerical dependencies in
\texttt{requirements.txt}. The article is self-contained and does not
require the repository's notation file. It can be compiled using
\begin{verbatim}
pdflatex -interaction=nonstopmode -halt-on-error article.tex
pdflatex -interaction=nonstopmode -halt-on-error article.tex
\end{verbatim}
The PDF was rendered and visually inspected during preparation. The
LaTeX source is the editable mathematical record; the numerical script is
supplementary evidence and an implementation of the displayed positive
sum, not a replacement for the proofs.

\begin{thebibliography}{9}
\bibitem{repo}
V.~Reshetnikov, \emph{ProveIt: Combinatorial Transseries and Their Inverses},
chapter ``The singular transition $q\to1$,'' source labels
\texttt{t3:prop:double-A}, \texttt{t3:eq:central-double-exact},
\texttt{t3:eq:central-double-asympt}, and \texttt{t3:eq:crossover-all}.
Repository snapshot 29 September 2026, commit
\texttt{04e06e032dff1966513bfba316966d52db3646b4}.
\url{https://github.com/VladimirReshetnikov/ProveIt/blob/04e06e032dff1966513bfba316966d52db3646b4/Analysis/FabiusFunction/docs/semi-formalized-research-frontiers/drafts/series-and-transseries/Combinatorial_Transseries_Inverses/Combinatorial_Transseries_Inverses.tex}.

\bibitem{reporeadme}
V.~Reshetnikov, \emph{ProveIt: Series and transseries}, group README and
\emph{Transseries and Inversion} package README, same snapshot.
\url{https://github.com/VladimirReshetnikov/ProveIt/blob/04e06e032dff1966513bfba316966d52db3646b4/Analysis/FabiusFunction/docs/semi-formalized-research-frontiers/drafts/series-and-transseries/README.md}.

\bibitem{costingaroufalidis}
O.~Costin and S.~Garoufalidis, Resurgence of the Euler--MacLaurin summation
formula, \emph{Annales de l'Institut Fourier} \textbf{58} (2008), 893--914.
DOI: \href{https://doi.org/10.5802/aif.2373}{10.5802/aif.2373}.
Preprint \url{https://arxiv.org/abs/math/0703641}.

\bibitem{garoufalidiskashaev}
S.~Garoufalidis and R.~Kashaev, \emph{Resurgence of Faddeev's quantum
dilogarithm}, arXiv:2008.12465v2, 5 October 2020.
\url{https://arxiv.org/abs/2008.12465v2}.

\bibitem{fantinirella}
V.~Fantini and C.~Rella, \emph{Modular resurgence, $q$-Pochhammer symbols,
and quantum operators from mirror curves}, arXiv:2506.08265v2,
1 April 2026. \url{https://arxiv.org/abs/2506.08265v2}.

\bibitem{dlmfcoth}
NIST Digital Library of Mathematical Functions, \S4.36,
Infinite Products and Partial Fractions for Hyperbolic Functions.
\url{https://dlmf.nist.gov/4.36}. Accessed 29 September 2026.

\bibitem{dlmfqgamma}
NIST Digital Library of Mathematical Functions, \S5.18,
$q$-Gamma and $q$-Beta Functions.
\url{https://dlmf.nist.gov/5.18}. Accessed 29 September 2026.

\bibitem{dlmfeta}
NIST Digital Library of Mathematical Functions, \S23.18,
Modular Transformations, including Dedekind's eta transformation.
\url{https://dlmf.nist.gov/23.18}. Accessed 29 September 2026.
\end{thebibliography}
\end{document}
